\documentclass[11pt,a4paper]{scrartcl}

\usepackage{amssymb}
\usepackage{amsmath}
\usepackage{systeme}

\usepackage[utf8]{inputenc}
\usepackage[T1]{fontenc}
\newcommand{\ats}[1]{\par\noindent\textit{Atsakymas:} #1}
\usepackage{cancel}
\usepackage{tikz}
\usepackage{lmodern} 
\usepackage{amsmath,amssymb,amsthm}
\usepackage{graphicx}
\usepackage[dvipsnames,svgnames]{xcolor}
\usepackage{hyperref}
\usepackage{mathrsfs}
\usepackage[shortlabels]{enumitem}
\usepackage[obeyFinal,textsize=scriptsize,shadow,loadshadowlibrary]{todonotes}
\usepackage{mathtools}
\usepackage{microtype}

%%% Paragraph layout
\setlength{\parindent}{0pt}
\setlength{\parskip}{0.6\baselineskip}

%%% Colored sections
\renewcommand*{\sectionformat}%
  {\color{purple}\S\thesection\enskip}
\renewcommand*{\subsectionformat}%
  {\color{purple}\S\thesubsection\enskip}
\renewcommand*{\subsubsectionformat}%
  {\color{purple}\S\thesubsubsection\enskip}
\addtokomafont{paragraph}{\color{orange!35!black}\P\ }
\KOMAoptions{numbers=noenddot}

%%% Title
\addtokomafont{subtitle}{\Large}
\setkomafont{author}{\Large\scshape}
\setkomafont{date}{\Large\normalsize}
% Neigiamas vertikalus tarpas, kuris pakelia dokumento antraštę į viršų. 
% Galite keisti -2.5cm į -3cm (aukščiau) arba -1.5cm (žemiau).
\titlehead{\vspace*{-7.5cm}} 

% PRIDĖTA: Automatiškai perrašome datą į mokyklos pavadinimą
\AtBeginDocument{%
  \date{\normalsize Vilniaus licėjus}
}

%%% Page setup
\usepackage[headsepline]{scrlayer-scrpage}
\renewcommand{\headfont}{}
\addtolength{\textheight}{3.14cm}
\setlength{\footskip}{0.5in}
\setlength{\headsep}{10pt}
% Puslapio antraštė turi rodyti tik konkrečius dokumento duomenis.
% Nustatome juos aiškiai, kad neatsirastų "\@author" / "\@title" arba senų reikšmių.
\makeatletter
\AtBeginDocument{%
  \ihead{\footnotesize\textbf{Andrius Berniukevičius}}%
  \ohead{\footnotesize\textbf{\@title}}%
}
\makeatother
\automark{section}
\chead{}
\cfoot{\pagemark}

%%% Macros
\providecommand{\ol}{\overline}
\providecommand{\ul}{\underline}
\providecommand{\wt}{\widetilde}
\providecommand{\wh}{\widehat}
\providecommand{\eps}{\varepsilon}
\providecommand{\half}{\frac{1}{2}}
\providecommand{\inv}{^{-1}}
\newcommand{\dang}{\measuredangle} %% Directed angle
\newcommand{\vocab}[1]{\textbf{\color{blue}\sffamily #1}}
\providecommand{\CC}{\mathbb C}
\providecommand{\FF}{\mathbb F}
\providecommand{\NN}{\mathbb N}
\providecommand{\QQ}{\mathbb Q}
\providecommand{\RR}{\mathbb R}
\providecommand{\ZZ}{\mathbb Z}
\providecommand{\ts}{\textsuperscript}
\providecommand{\dg}{^\circ}
\providecommand{\ii}{\item}
\providecommand{\defeq}{\coloneqq}
\DeclareMathOperator*{\lcm}{lcm}
\DeclareMathOperator*{\argmin}{arg min}
\DeclareMathOperator*{\argmax}{arg max}
\providecommand{\hrulebar}{\par
\hspace{\fill}\rule{0.95\linewidth}{.7pt}\hspace{\fill}
\par\nointerlineskip \vspace{\baselineskip}}

%%% Optional colored theorems
\usepackage{thmtools}
\usepackage[framemethod=TikZ]{mdframed}
\mdfdefinestyle{mdbluebox}{%
  roundcorner=10pt,
  linewidth=1pt,
  skipabove=12pt,
  innerbottommargin=9pt,
  skipbelow=2pt,
  linecolor=blue,
  nobreak=true,
  backgroundcolor=TealBlue!5,
}
\declaretheoremstyle[
  headfont=\sffamily\bfseries\color{MidnightBlue},
  mdframed={style=mdbluebox},
  headpunct={\\[3pt]},
  postheadspace={0pt}
]{thmbluebox}
\mdfdefinestyle{mdredbox}{%
  linewidth=0.5pt,
  skipabove=12pt,
  frametitleaboveskip=5pt,
  frametitlebelowskip=0pt,
  skipbelow=2pt,
  frametitlefont=\bfseries,
  innertopmargin=4pt,
  innerbottommargin=8pt,
  nobreak=true,
  backgroundcolor=Salmon!5,
  linecolor=RawSienna,
}
\declaretheoremstyle[
  headfont=\bfseries\color{RawSienna},
  mdframed={style=mdredbox},
  headpunct={\\[3pt]},
  postheadspace={0pt},
]{thmredbox}
\mdfdefinestyle{mdgreenbox}{%
  skipabove=8pt,
  linewidth=2pt,
  rightline=false,
  leftline=true,
  topline=false,
  bottomline=false,
  linecolor=ForestGreen,
  backgroundcolor=ForestGreen!5,
}
\declaretheoremstyle[
  headfont=\bfseries\sffamily\color{ForestGreen!70!black},
  bodyfont=\normalfont,
  spaceabove=2pt,
  spacebelow=1pt,
  mdframed={style=mdgreenbox},
  headpunct={ --- },
]{thmgreenbox}
\mdfdefinestyle{mdblackbox}{%
  skipabove=8pt,
  linewidth=3pt,
  rightline=false,
  leftline=true,
  topline=false,
  bottomline=false,
  linecolor=black,
  backgroundcolor=RedViolet!5!gray!5,
}
\declaretheoremstyle[
  headfont=\bfseries,
  bodyfont=\normalfont\small,
  spaceabove=0pt,
  spacebelow=0pt,
  mdframed={style=mdblackbox}
]{thmblackbox}

%% Aplinkos su rėmeliais (thmtools)
\declaretheorem[style=thmbluebox,name=Teorema]{theorem}
\declaretheorem[style=thmbluebox,name=Lema]{lemma}
\declaretheorem[style=thmbluebox,name=Savybė]{property}
\declaretheorem[style=thmbluebox,name=Teiginys]{proposition}
\declaretheorem[style=thmbluebox,name=Išvada]{corollary}
\declaretheorem[style=thmbluebox,name=Teorema,numbered=no]{theorem*}
\declaretheorem[style=thmbluebox,name=Lema,numbered=no]{lemma*}
\declaretheorem[style=thmbluebox,name=Teiginys,numbered=no]{proposition*}
\declaretheorem[style=thmbluebox,name=Išvada,numbered=no]{corollary*}
\declaretheorem[style=thmgreenbox,name=Algoritmas]{algorithm}
\declaretheorem[style=thmgreenbox,name=Algoritmas,numbered=no]{algorithm*}
\declaretheorem[style=thmgreenbox,name=Teiginys]{claim}
\declaretheorem[style=thmgreenbox,name=Teiginys,numbered=no]{claim*}
\declaretheorem[style=thmredbox,name=Pavyzdys]{example}
\declaretheorem[style=thmredbox,name=Pavyzdys,numbered=no]{example*}
\declaretheorem[style=thmblackbox,name=Pastaba]{remark}
\declaretheorem[style=thmblackbox,name=Pastaba,numbered=no]{remark*}

%% Standartinės amsthm aplinkos (Apibrėžimai ir užduotys)
\theoremstyle{definition}
\ifcsname conjecture\endcsname\else\newtheorem{conjecture}{Hipotezė}\fi
\ifcsname definition\endcsname\else\newtheorem{definition}{Apibrėžimas}\fi
\ifcsname fact\endcsname\else\newtheorem{fact}{Faktas}\fi
\ifcsname answer\endcsname\else\newtheorem{answer}{Atsakymas}\fi
\ifcsname ques\endcsname\else\newtheorem{ques}{Klausimas}\fi
\ifcsname exercise\endcsname\else\newtheorem{exercise}{Užduotis}\fi
\ifcsname problem\endcsname\else\newtheorem{problem}{Uždavinys}[section]\fi
\ifcsname conjecture*\endcsname\else\newtheorem*{conjecture*}{Hipotezė}\fi
\ifcsname definition*\endcsname\else\newtheorem*{definition*}{Apibrėžimas}\fi
\ifcsname fact*\endcsname\else\newtheorem*{fact*}{Faktas}\fi
\ifcsname answer*\endcsname\else\newtheorem*{answer*}{Atsakymas}\fi
\ifcsname ques*\endcsname\else\newtheorem*{ques*}{Klausimas}\fi
\ifcsname exercise*\endcsname\else\newtheorem*{exercise*}{Užduotis}\fi
\ifcsname problem*\endcsname\else\newtheorem*{problem*}{Uždavinys}\fi

\newcounter{answerssection}
\newcommand{\answerssection}{%
  \setcounter{answerssection}{\value{section}}%
  \section{Atsakymai}%
  \setcounter{problem}{0}%
  \renewcommand{\theproblem}{\arabic{answerssection}.\arabic{problem}}%
}

%% GALUTINIS SPRENDIMAS PROOF APLINKAI
%% --- PRADŽIA: Įrodymo aplinkos sutvarkymas ---
\makeatletter

% 1. Užtikriname, kad žodis būtų "Įrodymas", net jei "babel" paketas bando jį perrašyti
\AtBeginDocument{%
  \renewcommand{\proofname}{Įrodymas}%
  \@ifpackageloaded{babel}{%
    \addto\captionslithuanian{\renewcommand{\proofname}{Įrodymas}}%
  }{}%
}

% 2. Priverstinai pašaliname scrartcl klasės bold šriftą ir paliekame TIK italic
% 3. Sujungiame su tuščiu kvadratėliu dešiniajame kampe.
\renewcommand{\qedsymbol}{\ensuremath{\Box}}
\renewenvironment{proof}[1][\proofname]{\par
  \pushQED{\qed}%
  \normalfont \topsep6\p@\@plus6\p@\relax
  \trivlist
  \item[\hskip\labelsep
        \normalfont\itshape % Priverčia naudoti tik pasvirą šriftą, kaip nuotraukoje
    #1\@addpunct{.}]\ignorespaces
}{%
  \popQED\endtrivlist\@endpefalse
}
\newenvironment{solution}{\par
  \normalfont \topsep6\p@\@plus6\p@\relax
  \trivlist
  \item[\hskip\labelsep
        \normalfont\itshape
    \textit{Sprendimas}\@addpunct{.}]\ignorespaces
}{%
  \endtrivlist\@endpefalse
}
\newcommand{\ans}[1]{\par\noindent\textit{Atsakymas:} #1}
\makeatother


\title{Algebra I: Lygčių sistemos su trimis kintamaisiais}
\author{Andrius Berniukevičius}

\begin{document}

\maketitle

\section{Lygčių sistemos su trimis kintamaisiais}

Trijų kintamųjų lygčių sistemose jau nebeužtenka mechaniškai išreikšti vieną kintamąjį per kitus. Dažnai lemiama tampa pačios sistemos struktūra: sandaugos, cikliškumas, simetrija arba galimybė sudėti bei atimti lygtis.

Todėl tokią sistemą verta nagrinėti kaip vieną objektą, o ne kaip atskirų lygčių rinkinį. Tinkamai sujungus lygtis dažnai gaunamos bendros kintamųjų sandaugos, kvadratų sumos arba išskaidomosios išraiškos.

\section{Pagrindiniai metodai}

Sprendžiant sistemas su trimis kintamaisiais dažniausiai pasiteisina šie veiksmai:
\begin{itemize}
    \item sudėti lygtis, kad atsirastų kvadratų suma arba kita bendra išraiška;
    \item sudauginti lygtis, kad būtų galima rasti bendrą sandaugą, pavyzdžiui, $xyz$;
    \item atimti lygtis, kad išryškėtų dauginamieji ir kintamųjų ryšiai.
\end{itemize}

\subsection{Du naudingi faktai}
Vėlesniuose uždaviniuose prireiks dviejų faktų, kurie padeda atpažinti svarbią sistemos struktūrą. Pirmasis leidžia kubų sumą perrašyti per kintamųjų sumą, porines sandaugas ir sandaugą $xyz$:
$$
x^3+y^3+z^3=(x+y+z)(x^2+y^2+z^2-xy-xz-yz)+3xyz.
$$

Antrasis faktas leidžia iš kvadratų sumos, lygios nuliui, nustatyti kiekvieno dėmens reikšmę:

\begin{equation*}
\text{Jei }x_1^2+x_2^2+\dots+x_n^2=0 \text{, tai }x_1=x_2=\dots=x_n=0.
\end{equation*}


% ==========================================
% PAVYZDŽIAI
% ==========================================
\section{Uždavinių sprendimo pavyzdžiai}

\begin{example}[Sistemos daugyba]
Išspręskite lygčių sistemą:
\[
\left\{\begin{aligned}
    xy &= -30 \\
    yz &= -12 \\
    zx &= 10
\end{aligned}\right.
\]
\end{example}

\begin{solution}
Matome, kad visos lygtys yra grynos sandaugos. Tai klasikinis ženklas, kad geriausia strategija – sudauginti visas tris lygtis tarpusavyje. Prieš tai padarydami, pastebėkime svarbią detalę: kadangi visų sandaugų rezultatai nelygūs nuliui, tai nė vienas kintamasis negali būti nulis ($x \neq 0, y \neq 0, z \neq 0$).

Sudauginame visas kaires lygčių puses ir visas dešines puses:
$$ (xy)(yz)(zx) = (-30) \cdot (-12) \cdot 10 $$
Kairėje pusėje kiekviena raidė pasikartoja po du kartus, todėl sugrupuojame jas po vienu kvadratu:
$$ x^2 y^2 z^2 = 3600 $$
$$ (xyz)^2 = 60^2 $$

Ištraukę kvadratinę šaknį gauname du galimus bendrosios sandaugos variantus:
$$ xyz = 60 \quad \text{arba} \quad xyz = -60 $$

Dabar nagrinėjame pirmąjį atvejį: $xyz = 60$.
Norint rasti $x$, tereikia bendrą sandaugą padalinti iš antrosios lygties ($yz = -12$):
$$ x = \frac{xyz}{yz} = \frac{60}{-12} = -5 $$
Norint rasti $y$, dalijame bendrą sandaugą iš trečiosios lygties ($zx = 10$):
$$ y = \frac{xyz}{zx} = \frac{60}{10} = 6 $$
Norint rasti $z$, dalijame iš pirmosios lygties ($xy = -30$):
$$ z = \frac{xyz}{xy} = \frac{60}{-30} = -2 $$
Gavome pirmąjį atsakymų trejetą: $(x,y,z) = (-5, 6, -2)$.

Analogiškai nagrinėjame antrąjį atvejį: $xyz = -60$.
$$ x = \frac{-60}{-12} = 5, \quad y = \frac{-60}{10} = -6, \quad z = \frac{-60}{-30} = 2 $$
Gavome antrąjį trejetą: $(x,y,z) = (5, -6, 2)$. Abu trejetai tenkina pradinę sistemą.
\ats{$(-5; 6; -2)$ ir $(5; -6; 2)$}
\end{solution}


\begin{example}[Ciklinė suma]
Išspręskite lygčių sistemą:
\[
\left\{\begin{aligned}
    x^2 + y^2 &= 13 \\
    y^2 + z^2 &= 25 \\
    z^2 + x^2 &= 20
\end{aligned}\right.
\]
\end{example}

\begin{solution}
Išraiška arba lygčių sistema vadinama \vocab{cikline}, kurioje kiekvienas narys gaunamas ankstesniajame nuosekliai pakeitus $x \to y$, $y \to z$ ir $z \to x$. Šioje sistemoje pirmoji lygtis $x^2+y^2=13$ taip pereina į $y^2+z^2=25$, o ši -- į $z^2+x^2=20$.

Tokiose sistemose dažnai naudinga sudėti visas lygtis, nes kiekvienas kintamasis pasikartoja vienodai daug kartų. Išreikšti vieną kintamąjį ir statyti į kitą čia būtų neefektyvu.

Sudedame visas tris lygtis tarpusavyje. Kiekvienas kintamasis kairėje pusėje pasikartoja lygiai du kartus:
$$ 2x^2 + 2y^2 + 2z^2 = 13 + 25 + 20 $$
$$ 2(x^2 + y^2 + z^2) = 58 $$
Padalinę iš dviejų, gauname bazinę visų trijų kvadratų sumą:
$$ x^2 + y^2 + z^2 = 29 $$

Iš šios bendros sumos atimkime pirmąją pradinę lygtį ($x^2 + y^2 = 13$):
$$ (x^2 + y^2 + z^2) - (x^2 + y^2) = 29 - 13 $$
$$ z^2 = 16 \implies z = \pm 4 $$

Analogiškai atimame antrąją ir trečiąją lygtis iš bendros sumos:
$$ x^2 = 29 - 25 = 4 \implies x = \pm 2 $$
$$ y^2 = 29 - 20 = 9 \implies y = \pm 3 $$

Kadangi pradinėse lygybėse nėra sandaugų tarp skirtingų kintamųjų, kintamųjų ženklai vienas nuo kito nepriklauso. Vadinasi, galime laisvai kombinuoti pliusus ir minusus. Iš viso gauname $2 \cdot 2 \cdot 2 = 8$ skirtingus sprendinius.
\ats{$(\pm 2; \pm 3; \pm 4)$ (iš viso 8 sprendiniai)}
\end{solution}


\begin{example}[Sistemos sudėtis ir paslėpti kvadratai]
Raskite visus realiųjų skaičių trejetus $(x,y,z)$, tenkinančius sistemą:
\[
\left\{\begin{aligned}
    x^2+9 &= 4y \\
    y^2+1 &= 6z \\
    z^2+4 &= 2x
\end{aligned}\right.
\]
\end{example}

\begin{solution}
Šįkart lygtys sudarytos iš atskirų, pliusais sujungtų dėmenų, o laipsniai asimetriški. Joks tiesioginis įstatymas čia nepadės, todėl vėl ieškosime paslėptos struktūros sumuodami.

Sudedame visas tris lygtis (visą kairę ir visą dešinę pusę):
$$ x^2+y^2+z^2+14 = 2x+4y+6z $$

Viską sukeliame į kairę lygties pusę, prilygindami nuliui, ir sugrupuojame narius su ta pačia raide šalia vienas kito:
$$ (x^2-2x) + (y^2-4y) + (z^2-6z) + 14 = 0 $$

Skaičių 14 išskaidome taip, kad kiekvieniems skliaustams atitektų idealus skaičius, reikalingas pilnam kvadratui suformuoti.
Norint $x^2-2x$ paversti į $(x-1)^2$, trūksta $1$.
Norint $y^2-4y$ paversti į $(y-2)^2$, trūksta $4$.
Norint $z^2-6z$ paversti į $(z-3)^2$, trūksta $9$.
Kadangi $1 + 4 + 9 = 14$, mūsų turimas laisvasis narys pasidalina tobulai:
$$ (x^2-2x+1) + (y^2-4y+4) + (z^2-6z+9) = 0 $$
$$ (x-1)^2 + (y-2)^2 + (z-3)^2 = 0 $$

Pagal nulinių kvadratų taisyklę, neneigiamų skaičių suma lygi nuliui tik tokiu atveju, jeigu kiekvienas skaičius atskirai yra lygus nuliui:
$$ x-1=0 \implies x=1 $$
$$ y-2=0 \implies y=2 $$
$$ z-3=0 \implies z=3 $$
Atradome vienintelį potencialų kandidatą $(1, 2, 3)$.

Tai galimas sprendinys, gautas iš \textit{sudėtos} sistemos. Jį privaloma patikrinti pradinėse lygybėse.
Įstatome gautas reikšmes į pirmąją lygtį $x^2+9=4y$:
$$ 1^2 + 9 = 10, \quad \text{bet dešinėje pusėje} \quad 4 \cdot 2 = 8 $$
Dešimt nelygu aštuoniems. Lygybė negalioja! Vadinasi, mūsų rastas trejetas nėra pradinės sistemos sprendinys.
\ats{Sprendinių nėra.}
\end{solution}


\begin{example}[Lygčių atimtis ir simetrijos laužymas]
Raskite visus \textbf{teigiamus} realiuosius skaičius $x, y$ ir $z$, tenkinančius sistemą:
\[
\left\{\begin{aligned}
    x + y &= z^2 \\
    y + z &= x^2 \\
    z + x &= y^2
\end{aligned}\right.
\]
\end{example}

\begin{solution}
Ši sistema yra ciklinė: kiekviena kita lygtis gaunama pakeitus $x \to y$, $y \to z$ ir $z \to x$. Kaip ir ankstesniame pavyzdyje, pirmiausia galime sudėti visas lygtis:
$$ 2x+2y+2z=x^2+y^2+z^2. $$
Tačiau ši bendra lygybė dar neleidžia nustatyti atskirų kintamųjų reikšmių. Todėl, pasinaudodami cikliškumu, atimsime lygtis poromis.
Iš pirmosios lygties atimkime antrąją. Atimdami pastebime, kad kintamasis $y$ susiprastina:
$$ (x + y) - (y + z) = z^2 - x^2 $$
$$ x - z = z^2 - x^2 $$

Dešinėje pusėje matome kvadratų skirtumą. Išskaidykime jį dauginamaisiais:
$$ x - z = (z - x)(z + x) $$
Iškeliame minusą prieš skliaustus dešinėje, kad skliaustų turinys taptų vienodas su kairiąja puse:
$$ x - z = -(x - z)(x + z) $$
Sukeliame viską į kairę pusę ir iškeliame bendrą dauginamąjį $(x - z)$:
$$ (x - z) + (x - z)(x + z) = 0 $$
$$ (x - z)(1 + x + z) = 0 $$

Sandauga lygi nuliui tada, kai bent vienas iš dauginamųjų lygus nuliui.
Kadangi sąlygoje pasakyta, jog ieškome tik \textbf{teigiamų} skaičių, suma $(1 + x + z)$ privalo būti griežtai teigiama ($>0$). Vadinasi, nuliui lygiu būti gali tik pirmasis skliaustas:
$$ x - z = 0 \implies x = z $$

Analogiškai atėmę antrąją lygtį iš trečiosios, gautume $(y - x)(1 + y + x) = 0$, kas dėl tos pačios priežasties duoda $y = x$.
Padarome išvadą, kad visi trys kintamieji privalo būti lygūs: $x = y = z$.

Grįžtame į pradinę pirmąją lygtį ir visur vietoje $y$ bei $z$ įstatome $x$:
$$ x + x = x^2 $$
$$ x^2 - 2x = 0 $$
$$ x(x - 2) = 0 $$
Sprendiniai yra $x=0$ arba $x=2$. Kadangi ieškome tik teigiamų skaičių, $0$ atkrenta. Vadinasi, $x=2$, o kadangi visi kintamieji lygūs, tai ir $y=2$, ir $z=2$.
\ats{$(2, 2, 2)$}
\end{solution}


\vspace{0.5cm}

% ==========================================
% UŽDAVINIAI
% ==========================================
\newpage
\section{Uždaviniai savarankiškam sprendimui (14 iššūkių)}

\begin{enumerate}
    \renewcommand{\labelenumi}{\textbf{\theenumi.}}

    \item Išspręskite lygčių sistemą:
    \[
    \left\{
    \begin{array}{l}
    x+y=8\\
    y+z=10\\
    z+x=12
    \end{array}
    \right.
    \]

    \item Raskite teigiamus sprendinius:
    \[
    \left\{
    \begin{array}{l}
    xy=2\\
    yz=6\\
    zx=3
    \end{array}
    \right.
    \]

    \item Išspręskite lygčių sistemą:
    \[
    \left\{
    \begin{array}{l}
    x^2+y^2=10\\
    y^2+z^2=13\\
    z^2+x^2=5
    \end{array}
    \right.
    \]

    \item Raskite visus sprendinius:
    \[
    \left\{
    \begin{array}{l}
    a+b+c+d=4\\
    a+b+c+e=8\\
    a+b+d+e=12\\
    a+c+d+e=16\\
    b+c+d+e=20
    \end{array}
    \right.
    \]

    \item Išspręskite lygčių sistemą:
    \[
    \left\{
    \begin{array}{l}
    x(y+z)=5\\
    y(z+x)=8\\
    z(x+y)=9
    \end{array}
    \right.
    \]

    \item Raskite visus realiuosius skaičius $x$, $y$ ir $z$, kurie tenkina dvigubą lygybę:
    \[ x^2+y^2+z^2=x-z=2 \]

    \item Išspręskite lygčių sistemą:
    \[
    \left\{
    \begin{array}{l}
    2x^2-2x+2=y+z\\
    2y^2-2y+2=z+x\\
    2z^2-2z+2=x+y
    \end{array}
    \right.
    \]

    \item Raskite lygčių sistemos visus realiuosius sprendinius $(a,b,c,d)$:
    \[
    \begin{cases} 
    \begin{aligned}
    ab+c+d&=0 \\
    bc+d+a &= 8\\
    cd+a+b&=1\\
    da+b+c&=7
    \end{aligned}
    \end{cases}
    \]

    \item Išspręskite lygčių sistemą:
    \[
    \begin{cases}
    \begin{aligned} 
    x(y+z+1)&=y^2+z^2-5 \\
    y(z+x+1)&=z^2+x^2-5\\
    z(x+y+1)&=x^2+y^2-5
    \end{aligned}
    \end{cases}
    \]

    \item Raskite visus sprendinius:
    \[
    \left\{
    \begin{array}{l}
    \dfrac{x}{y}+\dfrac{y}{z}+\dfrac{z}{x}=3\\[0.3cm]
    \dfrac{y}{x}+\dfrac{z}{y}+\dfrac{x}{z}=3\\[0.3cm]
    \dfrac{1}{x}+\dfrac{1}{y}+\dfrac{1}{z}=2
    \end{array}
    \right.
    \]

    \item Duoti teigiami realieji skaičiai $x,y,z$, kurie tenkina sistemą:
    \[
    \left\{
    \begin{array}{l}
    x+\dfrac{1}{y}=4\\[0.2cm]
    y+\dfrac{1}{z}=1\\[0.2cm]
    z+\dfrac{1}{x}=\dfrac{7}{3}
    \end{array}
    \right.
    \]
    Raskite sandaugą $xyz$.

    \item Skaičiai $x$, $y$ ir $z$ sieja lygybės:
    $$x+\frac{1}{y}=y+\frac{1}{z}=z+\frac{1}{x}$$
    Įrodykite, kad arba $x=y=z$, arba $(xyz)^2=1$.

    \item Išspręskite lygčių sistemą:
    \[
    \left\{
    \begin{array}{l}
    x^2+4y^2-2z^4=18\\
    x+2y+2z^2=40\\
    z^4-2xy=9
    \end{array}
    \right.
    \]

    \item Raskite lygčių sistemos visus realiuosius sprendinius $(x,y,z)$:
    \[
    \left\{
    \begin{aligned}
    26x^2+5y^2+29z^2 &= 20xy+4xz+10yz,\\
    x^3+y^3+z^3 &= 1072
    \end{aligned}
    \right.
    \]
\end{enumerate}

% ==========================================
% ATSAKYMAI
% ==========================================
\newpage
\section{Atsakymai}

\begin{enumerate}
    \item $(5, 3, 7)$
    \item $(1, 2, 3)$
    \item $(\pm 1, \pm 3, \pm 2)$ (viso 8 sprendinių trejetai)
    \item $(-5, -1, 3, 7, 11)$
    \item $(1, 2, 3)$ ir $(-1, -2, -3)$
    \item $(1, 0, -1)$
    \item $(1, 1, 1)$
    \item $(1, 2, 3, 0), (1, -2, -3, 0), (-1, -2, 3, 0)$ ir $(-1, 2, -3, 0)$
    \item $(2, 2, 2)$ ir $(-2, -2, -2)$
    \item $\left(\frac{3}{2}, \frac{3}{2}, \frac{3}{2}\right)$
    \item $1$
    \item Įrodyta.
    \item $\left(13, \frac{7}{2}, \pm\sqrt{10}\right)$ ir $\left(7, \frac{13}{2}, \pm\sqrt{10}\right)$
    \item $(4, 10, 2)$
\end{enumerate}

\end{document}